\documentclass[12pt,a4paper]{article}
\usepackage[utf8]{inputenc}
\usepackage[T1]{fontenc}
\usepackage[spanish]{babel}
\usepackage{geometry}
\geometry{a4paper, top=2.5cm, bottom=2.5cm, left=2.5cm, right=2.5cm}
\usepackage{listings}
\usepackage{xcolor}
\usepackage{tcolorbox}
\usepackage{hyperref}
\usepackage{graphicx}
\usepackage{tikz}
\usetikzlibrary{positioning, arrows.meta}

\definecolor{codegreen}{rgb}{0,0.6,0}
\definecolor{codegray}{rgb}{0.5,0.5,0.5}
\definecolor{codepurple}{rgb}{0.58,0,0.82}
\definecolor{backcolour}{rgb}{0.96,0.96,0.96}

\lstdefinestyle{javastyle}{
    language=Java,
    backgroundcolor=\color{backcolour},
    commentstyle=\color{codegreen}\itshape,
    keywordstyle=\color{blue}\bfseries,
    numberstyle=\tiny\color{codegray},
    stringstyle=\color{codepurple},
    basicstyle=\ttfamily\footnotesize,
    breakatwhitespace=false,
    breaklines=true,
    captionpos=t,
    keepspaces=true,
    numbers=left,
    numbersep=5pt,
    showspaces=false,
    showstringspaces=false,
    showtabs=false,
    tabsize=4,
    frame=single,
    rulecolor=\color{lightgray}
}
\lstset{
    style=javastyle,
    literate={á}{{\'a}}1 {é}{{\'e}}1 {í}{{\'i}}1 {ó}{{\'o}}1 {ú}{{\'u}}1
             {Á}{{\'A}}1 {É}{{\'E}}1 {Í}{{\'I}}1 {Ó}{{\'O}}1 {Ú}{{\'U}}1
             {ñ}{{\~n}}1 {Ñ}{{\~N}}1
}

\title{\textbf{Guía de Solución: Programación Orientada a Objetos} \\ \large Actividad 3: Ejercicio 8.3 - Gestión de Eventos}
\author{}
\date{}

\begin{document}
\maketitle

\section{Repositorio de GitHub}
El código fuente de este proyecto se encuentra alojado en GitHub. 
\newline
\textbf{Enlace del repositorio:} \url{https://github.com/cioranemil/poo-trabajo-3-}

\section{Enunciado: Figuras Geométricas}
Se requiere desarrollar un programa con interfaz gráfica de usuario que permita calcular el volumen y superficie de varias figuras geométricas. Las figuras geométricas son el cilindro, la esfera y la pirámide.
\begin{itemize}
    \item Para el cilindro se solicitan su radio y altura.
    \item Para la esfera, su radio.
    \item Para la pirámide, su base, altura y apotema.
\end{itemize}
Una vez ingresados estos datos, el programa calcula el volumen y superficie de cada figura. Para desarrollar el programa se debe crear una jerarquía de clases para las diferentes figuras.

\section{Casos de Uso}
\begin{table}[h]
    \centering
    \begin{tabular}{|p{3cm}|p{12cm}|}
    \hline
    \textbf{Caso de Uso} & \textbf{Calcular Volumen y Superficie} \\ \hline
    \textbf{Actor} & Usuario / Estudiante \\ \hline
    \textbf{Descripción} & El usuario selecciona una figura tridimensional en el menú principal y realiza los cálculos de volumen y superficie ingresando sus medidas básicas. \\ \hline
    \textbf{Flujo Principal} & 
    1. El usuario visualiza el Menú Principal (\texttt{VentanaPrincipal}).\newline
    2. Hace clic en el botón de la figura deseada (Ej. Cilindro).\newline
    3. Se abre la ventana específica.\newline
    4. Introduce las dimensiones y presiona ``Calcular''.\newline
    5. El sistema procesa los datos y muestra el Volumen y Superficie. \\ \hline
    \end{tabular}
    \caption{Descripción del Caso de Uso}
\end{table}

\section{Diagrama de Clases y Explicación}

\begin{figure}[h]
    \centering
    \resizebox{\textwidth}{!}{
    \begin{tikzpicture}[
        node distance=1.5cm and 0.5cm,
        umlbox/.style={draw, inner sep=0pt, align=left, font=\footnotesize}
    ]
    
    \node[umlbox, align=center] (figura) {
        \begin{tabular}{l}
        \multicolumn{1}{c}{\textbf{FiguraGeometrica}} \\
        \hline
        - volumen: double \\
        - superficie: double \\
        \hline
        + setVolumen() \\
        + setSuperficie() \\
        + getVolumen() \\
        + getSuperficie()
        \end{tabular}
    };
    
    \node[umlbox, below=1.5cm of figura] (esfera) {
        \begin{tabular}{l}
        \multicolumn{1}{c}{\textbf{Esfera}} \\
        \hline
        - radio: double \\
        \hline
        + Esfera()
        \end{tabular}
    };
    
    \node[umlbox, left=0.5cm of esfera] (cilindro) {
        \begin{tabular}{l}
        \multicolumn{1}{c}{\textbf{Cilindro}} \\
        \hline
        - radio: double \\
        - altura: double \\
        \hline
        + Cilindro()
        \end{tabular}
    };
    
    \node[umlbox, right=0.5cm of esfera] (piramide) {
        \begin{tabular}{l}
        \multicolumn{1}{c}{\textbf{Piramide}} \\
        \hline
        - base: double \\
        - altura: double \\
        - apotema: double \\
        \hline
        + Piramide()
        \end{tabular}
    };
    
    \draw[-{Latex[open,length=4mm,width=4mm]}] (cilindro.north) -- +(0,0.5) -| (figura.south);
    \draw[-{Latex[open,length=4mm,width=4mm]}] (esfera.north) -- (figura.south);
    \draw[-{Latex[open,length=4mm,width=4mm]}] (piramide.north) -- +(0,0.5) -| (figura.south);
    
    \node[umlbox, below=1.5cm of esfera] (vesfera) {
        \begin{tabular}{l}
        \multicolumn{1}{c}{\textbf{VentanaEsfera}} \\
        \end{tabular}
    };
    \node[umlbox, below=1.5cm of cilindro] (vcilindro) {
        \begin{tabular}{l}
        \multicolumn{1}{c}{\textbf{VentanaCilindro}} \\
        \end{tabular}
    };
    \node[umlbox, below=1.5cm of piramide] (vpiramide) {
        \begin{tabular}{l}
        \multicolumn{1}{c}{\textbf{VentanaPiramide}} \\
        \end{tabular}
    };
    
    \draw[-] (vcilindro) -- node[left] {1} node[above, near end] {1} (cilindro);
    \draw[-] (vesfera) -- node[left] {1} node[above, near end] {1} (esfera);
    \draw[-] (vpiramide) -- node[left] {1} node[above, near end] {1} (piramide);
    
    \node[umlbox, below=1.5cm of vesfera] (vprincipal) {
        \begin{tabular}{l}
        \multicolumn{1}{c}{\textbf{VentanaPrincipal}} \\
        \end{tabular}
    };
    
    \draw[-] (vprincipal.north) -- +(0,0.5) -| node[above, near end] {1} (vcilindro.south);
    \draw[-] (vprincipal.north) -- node[left, near start] {1} node[left, near end] {1} (vesfera.south);
    \draw[-] (vprincipal.north) -- +(0,0.5) -| node[above, near end] {1} (vpiramide.south);
    
    \node[umlbox, below=1cm of vprincipal] (principal) {
        \begin{tabular}{l}
        \multicolumn{1}{c}{\textbf{Principal}} \\
        \end{tabular}
    };
    
    \draw[-] (principal) -- node[left] {1} node[left, near end] {1} (vprincipal);
    
    \end{tikzpicture}
    }
    \caption{Diagrama de clases del ejercicio 8.3}
\end{figure}

\textbf{Explicación del diagrama de clases:}
Se ha definido un paquete denominado ``Figuras'' que incluye un conjunto de clases. El punto de entrada al programa es la clase \texttt{Principal} que cuenta con el método \texttt{main}. La clase \texttt{Principal} está relacionada con la clase \texttt{VentanaPrincipal} que posee atributos privados para identificar tres botones que abrirán las ventanas respectivas. La clase \texttt{VentanaPrincipal} está conectada por medio de relaciones de asociación con las clases \texttt{VentanaCilindro}, \texttt{VentanaEsfera} y \texttt{VentanaPiramide}. 

Cada ventana específica se conecta con su clase lógica (ej. \texttt{VentanaCilindro} crea un único objeto \texttt{Cilindro}). Finalmente, las clases \texttt{Cilindro}, \texttt{Esfera} y \texttt{Piramide} son subclases de \texttt{FiguraGeometrica}, la cual es la clase padre que define atributos compartidos de volumen y superficie.

\section{Diagrama de Objetos}
\begin{verbatim}
[ miVentanaPrincipal: VentanaPrincipal ]
      |
      +---- [ cilindro: VentanaCilindro ] ---- [ cilindro: Cilindro ]
      |
      +---- [ esfera: VentanaEsfera ] -------- [ esfera: Esfera ]
      |
      +---- [ piramide: VentanaPiramide ] ---- [ piramide: Piramide ]
\end{verbatim}

\section{Ejecución del Programa}
A continuación se muestran los espacios para añadir las capturas de las interfaces gráficas.

\begin{figure}[h]
    \centering
    % \includegraphics[width=0.4\textwidth]{ventana_principal.png}
    \vspace{2cm} 
    \caption{a) Ventana Principal}
\end{figure}
\begin{figure}[h]
    \centering
    % \includegraphics[width=0.3\textwidth]{ventana_cilindro.png}
    \vspace{2cm} 
    \caption{b) Cilindro}
\end{figure}
\begin{figure}[h]
    \centering
    % \includegraphics[width=0.3\textwidth]{ventana_esfera.png}
    \vspace{2cm} 
    \caption{c) Esfera}
\end{figure}
\begin{figure}[h]
    \centering
    % \includegraphics[width=0.3\textwidth]{ventana_piramide.png}
    \vspace{2cm} 
    \caption{d) Pirámide}
\end{figure}

\newpage
\section{Solución: Código de la Interfaz de Usuario y Lógica}
A continuación se presenta todo el código desarrollado para la lógica jerárquica y las interfaces gráficas interconectadas.

\subsection*{Clase: FiguraGeometrica}
\begin{lstlisting}[caption={Clase FiguraGeometrica.java}]
package Figuras;

public class FiguraGeometrica {
    private double volumen;
    private double superficie;

    public void setVolumen(double volumen) { this.volumen = volumen; }
    public void setSuperficie(double superficie) { this.superficie = superficie; }
    public double getVolumen() { return this.volumen; }
    public double getSuperficie() { return this.superficie; }
}
\end{lstlisting}

\subsection*{Clase: Cilindro}
\begin{lstlisting}[caption={Clase Cilindro.java}]
package Figuras;

public class Cilindro extends FiguraGeometrica {
    private double radio;
    private double altura;

    public Cilindro(double radio, double altura) {
        this.radio = radio;
        this.altura = altura;
        setVolumen(Math.PI * altura * Math.pow(radio, 2.0));
        setSuperficie((2.0 * Math.PI * radio * altura) + (2.0 * Math.PI * Math.pow(radio, 2.0)));
    }
}
\end{lstlisting}

\subsection*{Clase: Esfera}
\begin{lstlisting}[caption={Clase Esfera.java}]
package Figuras;

public class Esfera extends FiguraGeometrica {
    private double radio;

    public Esfera(double radio) {
        this.radio = radio;
        setVolumen((4.0 / 3.0) * Math.PI * Math.pow(radio, 3.0));
        setSuperficie(4.0 * Math.PI * Math.pow(radio, 2.0));
    }
}
\end{lstlisting}

\subsection*{Clase: Piramide}
\begin{lstlisting}[caption={Clase Piramide.java}]
package Figuras;

public class Piramide extends FiguraGeometrica {
    private double base;
    private double altura;
    private double apotema;

    public Piramide(double base, double altura, double apotema) {
        this.base = base;
        this.altura = altura;
        this.apotema = apotema;
        setVolumen((Math.pow(base, 2.0) * altura) / 3.0);
        setSuperficie(Math.pow(base, 2.0) + (2.0 * base * apotema));
    }
}
\end{lstlisting}

\subsection*{Clase: VentanaCilindro}
\begin{lstlisting}[caption={Clase VentanaCilindro.java}]
package Figuras;

import java.awt.*;
import java.awt.event.ActionEvent;
import java.awt.event.ActionListener;
import javax.swing.*;

public class VentanaCilindro extends JFrame implements ActionListener {
    private Container contenedor;
    private JLabel radio, altura, volumen, superficie;
    private JTextField campoRadio, campoAltura;
    private JButton calcular;

    public VentanaCilindro() {
        contenedor = getContentPane();
        contenedor.setLayout(null);
        setTitle("Cilindro");
        setSize(280, 210);
        setLocationRelativeTo(null);
        setResizable(false);

        radio = new JLabel("Radio (cms):"); radio.setBounds(20, 20, 135, 23);
        campoRadio = new JTextField(); campoRadio.setBounds(100, 20, 135, 23);

        altura = new JLabel("Altura (cms):"); altura.setBounds(20, 50, 135, 23);
        campoAltura = new JTextField(); campoAltura.setBounds(100, 50, 135, 23);

        calcular = new JButton("Calcular"); calcular.setBounds(100, 80, 135, 23);
        calcular.addActionListener(this);

        volumen = new JLabel("Volumen (cm3):"); volumen.setBounds(20, 110, 230, 23);
        superficie = new JLabel("Superficie (cm2):"); superficie.setBounds(20, 140, 230, 23);

        contenedor.add(radio); contenedor.add(campoRadio);
        contenedor.add(altura); contenedor.add(campoAltura);
        contenedor.add(calcular); contenedor.add(volumen); contenedor.add(superficie);
    }

    @Override
    public void actionPerformed(ActionEvent event) {
        try {
            double r = Double.parseDouble(campoRadio.getText());
            double a = Double.parseDouble(campoAltura.getText());
            Cilindro cilindro = new Cilindro(r, a);
            volumen.setText("Volumen (cm3): " + String.format("%.2f", cilindro.getVolumen()));
            superficie.setText("Superficie (cm2): " + String.format("%.2f", cilindro.getSuperficie()));
        } catch (Exception e) {
            JOptionPane.showMessageDialog(null, "Campo nulo o error de formato",
                                          "Error", JOptionPane.ERROR_MESSAGE);
        }
    }
}
\end{lstlisting}

\subsection*{Clase: VentanaEsfera}
\begin{lstlisting}[caption={Clase VentanaEsfera.java}]
package Figuras;

import java.awt.*;
import java.awt.event.ActionEvent;
import java.awt.event.ActionListener;
import javax.swing.*;

public class VentanaEsfera extends JFrame implements ActionListener {
    private Container contenedor;
    private JLabel radio, volumen, superficie;
    private JTextField campoRadio;
    private JButton calcular;

    public VentanaEsfera() {
        contenedor = getContentPane();
        contenedor.setLayout(null);
        setTitle("Esfera");
        setSize(280, 200);
        setLocationRelativeTo(null);
        setResizable(false);

        radio = new JLabel("Radio (cms):"); radio.setBounds(20, 20, 135, 23);
        campoRadio = new JTextField(); campoRadio.setBounds(100, 20, 135, 23);

        calcular = new JButton("Calcular"); calcular.setBounds(100, 50, 135, 23);
        calcular.addActionListener(this);

        volumen = new JLabel("Volumen (cm3):"); volumen.setBounds(20, 90, 230, 23);
        superficie = new JLabel("Superficie (cm2):"); superficie.setBounds(20, 120, 230, 23);

        contenedor.add(radio); contenedor.add(campoRadio);
        contenedor.add(calcular); contenedor.add(volumen); contenedor.add(superficie);
    }

    @Override
    public void actionPerformed(ActionEvent evento) {
        try {
            double r = Double.parseDouble(campoRadio.getText());
            Esfera esfera = new Esfera(r);
            volumen.setText("Volumen (cm3): " + String.format("%.2f", esfera.getVolumen()));
            superficie.setText("Superficie (cm2): " + String.format("%.2f", esfera.getSuperficie()));
        } catch (Exception e) {
            JOptionPane.showMessageDialog(null, "Campo nulo o error de formato",
                                          "Error", JOptionPane.ERROR_MESSAGE);
        }
    }
}
\end{lstlisting}

\subsection*{Clase: VentanaPiramide}
\begin{lstlisting}[caption={Clase VentanaPiramide.java}]
package Figuras;

import java.awt.*;
import java.awt.event.ActionEvent;
import java.awt.event.ActionListener;
import javax.swing.*;

public class VentanaPiramide extends JFrame implements ActionListener {
    private Container contenedor;
    private JLabel base, altura, apotema, volumen, superficie;
    private JTextField campoBase, campoAltura, campoApotema;
    private JButton calcular;

    public VentanaPiramide() {
        contenedor = getContentPane();
        contenedor.setLayout(null);
        setTitle("Piramide");
        setSize(280, 240);
        setLocationRelativeTo(null);
        setResizable(false);

        base = new JLabel("Base (cms):"); base.setBounds(20, 20, 135, 23);
        campoBase = new JTextField(); campoBase.setBounds(120, 20, 135, 23);

        altura = new JLabel("Altura (cms):"); altura.setBounds(20, 50, 135, 23);
        campoAltura = new JTextField(); campoAltura.setBounds(120, 50, 135, 23);

        apotema = new JLabel("Apotema (cms):"); apotema.setBounds(20, 80, 135, 23);
        campoApotema = new JTextField(); campoApotema.setBounds(120, 80, 135, 23);

        calcular = new JButton("Calcular"); calcular.setBounds(120, 110, 135, 23);
        calcular.addActionListener(this);

        volumen = new JLabel("Volumen (cm3):"); volumen.setBounds(20, 140, 230, 23);
        superficie = new JLabel("Superficie (cm2):"); superficie.setBounds(20, 170, 230, 23);

        contenedor.add(base); contenedor.add(campoBase);
        contenedor.add(altura); contenedor.add(campoAltura);
        contenedor.add(apotema); contenedor.add(campoApotema);
        contenedor.add(calcular); contenedor.add(volumen); contenedor.add(superficie);
    }

    @Override
    public void actionPerformed(ActionEvent event) {
        try {
            double b = Double.parseDouble(campoBase.getText());
            double a = Double.parseDouble(campoAltura.getText());
            double ap = Double.parseDouble(campoApotema.getText());
            Piramide piramide = new Piramide(b, a, ap);
            volumen.setText("Volumen (cm3): " + String.format("%.2f", piramide.getVolumen()));
            superficie.setText("Superficie (cm2): " + String.format("%.2f", piramide.getSuperficie()));
        } catch (Exception e) {
            JOptionPane.showMessageDialog(null, "Campo nulo o error de formato",
                                          "Error", JOptionPane.ERROR_MESSAGE);
        }
    }
}
\end{lstlisting}

\subsection*{Clase: VentanaPrincipal}
\begin{lstlisting}[caption={Clase VentanaPrincipal.java}]
package Figuras;

import java.awt.*;
import java.awt.event.ActionEvent;
import java.awt.event.ActionListener;
import javax.swing.*;

public class VentanaPrincipal extends JFrame implements ActionListener {
    private Container contenedor;
    private JButton cilindro, esfera, piramide;

    public VentanaPrincipal() {
        contenedor = getContentPane();
        contenedor.setLayout(null);
        setTitle("Figuras Geometricas");
        setSize(350, 160);
        setLocationRelativeTo(null);
        setDefaultCloseOperation(JFrame.EXIT_ON_CLOSE);

        cilindro = new JButton("Cilindro"); cilindro.setBounds(20, 50, 90, 23);
        cilindro.addActionListener(this);

        esfera = new JButton("Esfera"); esfera.setBounds(130, 50, 90, 23);
        esfera.addActionListener(this);

        piramide = new JButton("Piramide"); piramide.setBounds(240, 50, 90, 23);
        piramide.addActionListener(this);

        contenedor.add(cilindro); contenedor.add(esfera); contenedor.add(piramide);
    }

    @Override
    public void actionPerformed(ActionEvent evento) {
        if (evento.getSource() == esfera) {
            new VentanaEsfera().setVisible(true);
        }
        if (evento.getSource() == cilindro) {
            new VentanaCilindro().setVisible(true);
        }
        if (evento.getSource() == piramide) {
            new VentanaPiramide().setVisible(true);
        }
    }
}
\end{lstlisting}

\subsection*{Clase: Principal}
\begin{lstlisting}[caption={Clase Principal.java}]
package Figuras;

public class Principal {
    public static void main(String[] args) {
        VentanaPrincipal menu = new VentanaPrincipal();
        menu.setVisible(true);
        menu.setResizable(false);
    }
}
\end{lstlisting}

\end{document}
